%% ===================================================================
\section{Introduction}
\label{sec:intro}
%% ===================================================================

Graph neural networks have entered domains where incorrect predictions carry tangible consequences: fraudulent accounts evade detection when an adversary inserts edges into a financial transaction graph~\cite{zugner2018adversarial}, drug interaction models produce unsafe recommendations when molecular graphs are perturbed~\cite{dai2018adversarial}, and power grid monitoring systems may miss critical contingencies if their graph-based representations are structurally fragile~\cite{nakiganda2023graph}. In all these settings, the graph topology is not merely an input format but a potential attack surface. A practitioner preparing to deploy a GNN faces a fundamental question that current tools leave unanswered: \emph{which edges in this graph, if perturbed, would cause the model's predictions to fail, and by how much?}

This question sits at the intersection of adversarial attack, certified defense, and sensitivity analysis, but no existing tool answers it completely. Attack methods such as Nettack~\cite{zugner2018adversarial} and Mettack~\cite{zugner2019adversarial} provide per-edge gradient information and per-node targeting, but optimize surrogate objectives (classification loss on a substitute model) rather than directly characterizing the sensitivity geometry, and they do not yield per-node tolerance budgets. Certified defenses based on randomized smoothing~\cite{bojchevski2020efficient,scholten2022randomized} or interval bound propagation~\cite{zugner2019certifiable,bojchevski2019certifiable,wang2021certified} provide robustness guarantees---including per-node certificates in the localized variant~\cite{schuchardt2023localized}---but do not identify \emph{which edges} are responsible for a node's vulnerability or provide the globally most sensitive perturbation direction. Empirical defenses~\cite{zhu2019robust,zhang2020gnnguard,jin2020graph,luo2021learning} harden the model but offer no formal understanding of residual vulnerability. What is missing is a unified \emph{structural vulnerability analysis} that, from a single computation under realistic perturbation constraints, yields the maximally sensitive perturbation direction, per-edge vulnerability rankings, and per-node tolerance budgets.

We introduce \AEGIS (Adversarial Evaluation of Graph Integrity via Sensitivity), a framework that produces exactly this map. The central object is the \emph{constrained sensitivity matrix} $S_c \in \R^{Nd \times |E|}$, which projects the full $N^2$-dimensional adjacency perturbation space onto the $|E|$-dimensional space of realistic (symmetric, edge-only) perturbations. From $S_c$, three outputs follow directly: (i)~the SVD yields the first-order maximally sensitive perturbation direction (optimal for equilibrium shift; empirically correlated with classification damage, \cref{sec:adaptive}), (ii)~column norms give per-edge vulnerability scores, and (iii)~the classification margin divided by the sensitivity norm gives per-node first-order tolerance radii (diagnostic quantities, not probabilistic certificates). The $S_c$ computation applies to GNNs with continuous edge-weight-modulated message passing (\cref{prop:explicit}; standard GAT and architectures using max/min aggregations require modification for continuous differentiability), yielding vulnerability rankings, SVD-optimal attacks, and first-order sensitivity radii as a practical computational tool. For contractive implicit GNNs (IGNN-class~\cite{gu2020implicit}), the implicit function theorem additionally provides formal guarantees: a critical perturbation budget $\ecrit = (1-\kappa)/\norm{W}_2$ and a three-regime characterization of model behavior under increasing perturbation (\cref{thm:phase_transition}). These formal guarantees are specific to the implicit case and do not transfer to explicit architectures.

These predictions are quantitatively accurate: constrained first-order predictions remain within 15\% of ground truth at $\varepsilon = 0.10$ across 7~GNN architectures and 9~datasets spanning citation networks, e-commerce, encyclopedias, and power grids. Continuous-to-discrete vulnerability transfer is positive in 29 of 33 architecture-dataset combinations, with strongest transfer for deeper models on sparse graphs (Kendall $\tau$ up to $+0.89$); practitioners should verify transfer for their specific architecture-graph pairing (\cref{sec:explicit_extension}). On IEEE power grids, the vulnerability spectrum correlates with N-1 contingency rankings (P@10 = 0.66--0.87), demonstrating that structural sensitivity analysis can serve as a pre-deployment screening layer for real-world engineering problems.

\textbf{Contributions.}
\begin{enumerate}
\item \textbf{Constrained sensitivity framework with architecture-specific guarantees.} We derive the constrained sensitivity matrix $S_c$, applicable as a \emph{computational tool} to any GNN with continuous edge-weight-modulated message passing (\cref{prop:explicit}; standard GAT and max/min aggregations require modification). From a single $S_c$ computation, three outputs follow: the first-order maximally sensitive perturbation direction, per-edge vulnerability rankings, and per-node first-order tolerance radii. For contractive implicit models (IGNN-class), we additionally derive a conservative sufficient condition on the critical budget $\ecrit$ and a three-regime vulnerability characterization (\cref{thm:phase_transition}); these \emph{formal guarantees} are specific to the implicit case and do not transfer to explicit architectures, which receive the computational analysis but not the convergence certificates.

\item \textbf{Scalable matrix-free computation.} A Neumann-series resolvent combined with autograd JVPs and randomized SVD~\cite{halko2011finding} computes $S_c$'s action without materializing any $Nd \times N^2$ matrix, enabling full-graph analysis (Cora $N{=}2{,}708$ in 78\,s). A comprehensive attack evaluation spanning gradient-based and gradient-free, continuous and discrete, same-objective and independent baselines (\cref{sec:adaptive}) confirms that the SVD direction is well-separated from random alternatives (best-of-1{,}000 random directions reaches only 48\% of $\sigma_1$) and validates rankings across 7 architectures (\cref{sec:explicit_extension}).

\item \textbf{Cross-domain validation and case study.} We validate across 9 datasets in 4 domains with 7 architectures (330 runs total). Continuous-to-discrete transfer is positive in 88\% of settings (Kendall $\tau$ from $-0.28$ to $+0.89$). A power grid case study demonstrates that $S_c$ vulnerability scores correlate with N-1 contingency rankings ($\tau = +0.37$--$+0.72$, P@10 $= 0.66$--$0.87$; \cref{sec:power_flow}), serving as a proof-of-concept for cross-domain transfer.
\end{enumerate}
